\begin{center}
  \large\textbf{ABSTRAK}
\end{center}

\addcontentsline{toc}{chapter}{ABSTRAK}

\vspace{2ex}

\begingroup
% Menghilangkan padding
\setlength{\tabcolsep}{0pt}

\noindent
\begin{tabularx}{\textwidth}{l >{\centering}m{2em} X}
  Nama Mahasiswa    & : & \name{}       \\

  Judul Tugas Akhir & : & \tatitle{}    \\

  Pembimbing        & : & 1. \advisor{} \\ & & 2. \coadvisor{} \\
\end{tabularx}
\endgroup

% Ubah paragraf berikut dengan abstrak dari tugas akhir
Sistem tiket elektronik memberikan kemudahan dalam pembelian dan manajemen
tiket serta mengurangi risiko kehilangan. Namun, tantangan akibat sentralisasi
masih menjadi kendala, termasuk kerentanan terhadap serangan siber dan
ketergantungan pada pihak ketiga. Teknologi blockchain hadir sebagai solusi
dengan menawarkan pencatatan transaksi yang terdesentralisasi, transparan, dan
aman. Teknologi ini juga mendukung inovasi aset digital seperti NFT, yang
membuka peluang baru dalam transaksi dan koleksi digital. Untuk mengatasi
keterbatasan biaya dan skalabilitas jaringan utama Ethereum, Polygon zkEVM
sebagai solusi Layer-2 berbasis rollup kriptografi memungkinkan transaksi yang
lebih efisien, cepat, dan murah tanpa mengorbankan keamanan. Meskipun masih
terdapat tantangan seperti perlindungan hak kekayaan intelektual dan
interoperabilitas antar-rantai, adopsi blockchain dan Polygon menunjukkan
potensi besar dalam membangun ekosistem aset digital yang inklusif dan efisien.
% Ubah kata-kata berikut dengan kata kunci dari tugas akhir
Kata Kunci: \emph{Blockchain}, Ticket Elektronik, \emph{NFT}, \emph{Polygon
  zkEVM}, Desentralisasi.
